\section{Introduction}
\label{sec:intro}

\subsection*{a. Objectives and research plan}
This paper evaluates a causal \emph{nowcast} of the reported carbon-impurity CES temperature on one pedestal-top chord of the KSTAR 2022 campaign; the data provider selected the chord. At each frame of a 10\,ms grid the estimator receives two kinds of input: CES reports that became available strictly before the frame, and current and past BES and ECEI voltage snapshots. From these it estimates the CES report belonging to the current frame before that report is available.

The work proceeds in two stages. In the first, history-only and lightweight models are compared with reference estimators on matched rows, information sets and computation budgets. The aim is to test whether other diagnostics add useful temporal information to the CES history. In the second, the measured errors, modality contributions and costs from that comparison are used to design a neural network for the task. This second stage is a plan; it is not reported here as a result.

Three comparisons organize the paper:
\begin{enumerate}\setlength{\itemsep}{1pt}\setlength{\parskip}{0pt}
\item Which causal estimators are useful when scored against three references? The references are hold, a causal Gaussian-process estimator, and two-sided interpolation of the monotone piecewise-cubic family~\citep{fritsch1980}. The last needs the next report, so it serves only as an offline reference (\S\ref{sec:res:best}).
\item What information do BES and ECEI add to the CES history (\S\ref{sec:res:fast})?
\item Do the answers persist under four validation designs, which are given equal weight, and does the computation fit within a 10\,ms frame (\S\ref{sec:res:designs})?
\end{enumerate}
Optional application analyses (\S\ref{sec:res:use}) and the conditions under which the estimator does not help (\S\ref{sec:res:limits}) follow.

Two targets are scored. F is the CES report of the current frame. L5 is a level computed offline as the median of at least three of the five subsequent reports. Both are targets only: neither the current F nor the future-derived L5 is an online input.

\subsection*{b. Why the problem is worth studying}
The edge ion temperature of H-mode plasmas changes on several time scales. On KSTAR, CES has shown the $T_i$ pedestal building up over tens to hundreds of milliseconds after the L--H transition and falling during ELM suppression by resonant magnetic perturbations~\citep{Ko2015IonTemperature}. On ASDEX Upgrade, a dedicated fast system recorded edge impurity temperature collapsing and recovering within a few milliseconds of type-I ELM crashes~\citep{Viezzer2018IonHeat}. Such intra-ELM dynamics lie below the 10\,ms grid used here.

Requirement studies for burning-plasma devices ask for comparable time resolution. Specifications developed for FIRE call for edge $T_i$ profiles at 10\,ms and about 10\,\% precision~\citep{Young2008PlasmaMeasurements}, and the ITER specifications list edge $T_i$ at 100\,ms and 10\,\%~\citep{Donne2007ChapterDiagnostics}. These are design targets for other devices, not KSTAR requirements. Charge-exchange $T_i$ has also served directly in feedback: on JT-60U, neutral-beam power regulated a two-chord core $T_i$ difference within a 10\,ms control cycle~\citep{Yoshida2009RealTime}.

CES infers impurity temperature and rotation from the Doppler width and shift of a line excited by charge exchange with beam neutrals~\citep{Fonck1984DeterminationPlasma}. The KSTAR instrument observes C\,VI near 529\,nm and has a separate background view~\citep{Ko2010ChargeExchange}. In the beam-modulation analysis described for KSTAR, the background of a beam-on frame is interpolated from the beam-off frames that bracket it~\citep{Lee2020TwoGaussian,Lee2026NeuralNetworkbased}. The authors of the former argue that this analysis cannot begin before the following beam-off frame. That is an argument about the procedure, not a measurement. We did not record when our reports became available, and we assign them no delay.

The reports are also conventional labels rather than independent physical truth. Active and passive emission must be separated in the analysis~\citep{Viezzer2011PassiveEdge}. On DIII-D, main-ion and carbon temperatures differed in the edge, by tens of eV inside the H-mode pedestal top and by more near the separatrix~\citep{Haskey2018MainIon}. Occasional fit failures have been associated with abrupt outliers in KSTAR CES profiles~\citep{Kim2024KineticProfile}. The simplest nowcast holds the last report, and this is the reference against which any added information must be shown.

\subsection*{c. Expected outcome}
The intended product is an estimate for frames whose CES report is not yet available, together with an assessment of its error. The assessment reports matched-row error differences against the references, the dependence on the age and availability of the CES history, the contribution of each modality, the sensitivity to validation design, and the computation time. Together these indicate when such an estimator could be used.

The estimator is complementary to spectrum reading. An upstream reader such as NN-CES converts a recorded spectrum into a report~\citep{Lee2026NeuralNetworkbased}. Our downstream estimator works from the report history and needs no current CES spectrum. A faster upstream reader would refresh the anchors on which the nowcast rests, and it could equally shrink the interval in which a nowcast is needed. Which effect prevails must be measured. We did not implement this coupling.

Meeting a computation budget is not a real-time deployment, and no physical mechanism is identified. The results concern one chord in one reused campaign, with no untouched data.

\subsection*{d. Relation to prior work}
Neural networks trained on conventionally fitted spectra have long been used to read charge-exchange line parameters. At JET, networks gave approximate line parameters~\citep{Bishop1993AutomaticAnalysis}, and a later monitoring demonstration produced $T_i$ profiles at the 50\,ms exposure cadence~\citep{Svensson1999AnalysisJet}. A recent example comes from HL-3~\citep{Tian2026RealTimeCxrs}. At DIII-D, a simplified real-time fit is reported to deliver values every 5\,ms with a delay of a few milliseconds; its agreement with the offline analysis is lower on edge chords than on core chords~\citep{Chrystal2026RealTime}. On KSTAR, a two-Gaussian fit~\citep{Lee2020TwoGaussian} and NN-CES~\citep{Lee2026NeuralNetworkbased} both dispense with beam-off frames. The NN-CES authors report close agreement with the modulation analysis inside $R\approx2.2$\,m and weaker agreement farther out. We have not examined NN-CES on our chord and claim no comparison with it.

Our task differs from forecasting future profiles. Neural-network forecasting has, for example, predicted profile changes about 200\,ms ahead from present profiles and proposed actuator trajectories~\citep{Abbate2021DatadrivenProfile}. A further line of work estimates missing or slow measurements from other signals. A DIII-D real-time kinetic reconstruction, trained with input dropout, still returns profiles, including an ion temperature profile, when charge-exchange data are absent; that ion temperature was assessed qualitatively~\citep{Shousha2023MachineLearningbased}. Faster diagnostics have been mapped to Thomson-scattering electron profiles at DIII-D~\citep{Jalalvand2025MultimodalSuperresolution} and COMPASS~\citep{Imrek2026TemporalSuperresolution}. On TCV, a Kalman observer merges fast interferometry with slower Thomson scattering~\citep{Pastore2023ModelBased}. We do not suggest that this literature is noncausal or restricted to electron quantities.

The present question is timed differently. It concerns one ion-temperature chord, with past reports of the same diagnostic as anchors. Its fast inputs are BES~\citep{Lampert2015CombinedHydrogen} and ECEI~\citep{Yun2014Quasi3d} voltage snapshots. Edge ECEI can respond to density as well as temperature~\citep{Vanovac2018EffectsDensity}, so the snapshots are treated as uncalibrated predictors, not as calibrated $T_e$ or $n_e$ and not as full-rate turbulence measurements. The estimators are built from familiar models, and no algorithmic priority is claimed. The contribution is an application-specific empirical evaluation.
